% Ad Familiares 12.30 (ad-familiares-12-30)
% Source: https://raw.githubusercontent.com/PerseusDL/canonical-latinLit/master/data/phi0474/phi056/phi0474.phi056.perseus-lat1.xml
% Fetched by scripts/fetch_latin.py

% Dateline (Perseus): Scr. Romae paulo post v Id. Iun. a. 711 (43).

\ciceroLetterOpener{CICERO CORNIFICIO S.}

\ciceroSection{1} itane ? praeter litigatores nemo ad te meas litteras? multae istae quidem ; tu enim perfecisti ut nemo sine litteris meis tibi se commendatum putaret ; sed quis umquam tuorum mihi dixit esse cui darem, quin dederim? aut quid mi iucundius quam, cum coram tecum loqui non possim, aut scribere ad te aut tuas legere litteras? illud magis mihi solet esse molestum, tantis me impediri occupationibus, ut ad te scribendi meo arbitratu facultas nulla detur. non enim te epistulis sed voluminibus lacesserem ; quibus quidem me a te provocari oportebat. quamvis enim occupatus sis,' oti tamen plus habes ; aut, si ne tu quidem vacas, noli impudens esse nec mihi molestiam exhibere et a me litteras crebriores, cum tu mihi raro mittas, flagitare.

\ciceroSection{2} nam cum antea distinebar maximis occupationibus, propterea quod omnibus curis rem p. mihi tuendam putabam, tum hoc tempore multo distineor vehementius. ut enim gravius aegrotant ii qui, cum levati morbo videntur, in eum de integro inciderunt, sic vehementius nos laboramus, qui profligato bello ac paene sublato renovatum bellum gerere conamur. sed haec hactenus. .

\ciceroSection{3} tu tibi, mi Cornifici, fac ut persuadeas non esse me tam imbecillo animo, ne dicam inhumano, ut a te vinci possim aut officiis aut amore. non dubitabam equidem, verum tamen multo mihi notiorem amorem tuum effecit Chaerippus. O hominem semper illum quidem mihi aptum, nunc vero etiam suavem! vultus me hercule tuos mihi expressit omnis, non solum animum ac verba pertulit. itaque noli vereri ne tibi suscensuerim, quod eodem exemplo ad me quo ad ceteros. requisivi equidem proprias ad me unum litteras, sed neque vehementer et amanter.

\ciceroSection{4} de sumptu, quem te in rem militarem facere et fecisse dicis, nihil sane possum tibi opitulari, propterea quod et orbus senatus consulibus amissis et incredibiles angustiae pecuniae publicae ; quae conquiritur undique, ut optime meritis militibus promissa solvantur ; quod quidem fieri sine tributo posse non arbitror.

\ciceroSection{5} de Attio Dionysio nihil puto esse, quoniam mihi nihil dixit Tratorius. de P. Lucceio nihil tibi concedo, quo studiosior tu sis quam ego sum ; est enim nobis necessarius. sed a magistris cum contenderem de praeterendo die, probarunt mihi sese quo minus id facerent et compromisso et iure iurando impediri. qua re veniendum arbitror Lucceio ; quamquam, si meis litteris obtemperavit, cum tu haec leges illum Romae esse oportebit.

\ciceroSection{6} ceteris de rebus maximeque de pecunia, cum Pansae mortem ignorares, scripsisti quae per nos ab eo consequi te posse arbitrarere. quae te non fefellissent, si viveret ; nam te diligebat ; post mortem autem eius quid fieri posset non videbamus. de Venuleio, Latino, Horatio valde laudo ; illud non nimium probo, quod scribis, quo illi animo aequiore ferrent, te tuis etiam legatis lictores ademisse (honore enim cum ignominia dignis non erant comparandi), eosque, ex s. c. si non decedunt, cogendos ut decedant existimo. haec fere ad eas litteras, quas eodem exemplo is binas accepi ; de reliquo velim tibi persuadeas non esse mihi meam dignitatem tua cariorem.
